% OWNER: theory-a
% STATUS: review (poprawiono: terminologia, usunięto potoczne Tinder na rzecz afektywnej selekcji kart, usunięto disclaimer ComfyUI, uzupełniono źródła)
% SOURCES: notes/SPIS_TRESCI_SZCZEGOLOWY.md §3; GLOSSARY.md; notes/HIPOTEZY_REWIZJA.md
\chapter{Preferencje estetyczne: jawne, ukryte i behawioralne}
\label{ch03:top}
\ChapterStatus{review}{theory-a}

Rozdział przedstawia strukturę preferencji estetycznych jako wieloźródłowe tworzywo procesu personalizacji w platformie IDA.
Wychodzi od klasycznego rozróżnienia na preferencje \textbf{jawne} (\emph{explicit}), deklarowane świadomie przez użytkownika, preferencje \textbf{ukryte} (\emph{implicit}), rejestrowane poprzez spontaniczne reakcje afektywne, oraz wskaźniki \textbf{behawioralne}. Profil osobowości (Big Five) stanowi trzeci filar modelu, szczegółowo rozwinięty w rozdziale~\ref{ch05:top}.
Celem metodycznym jest synteza wieloźródłowa, zasilająca eksperymentalną, ślepą macierz sześciu wariantów koncepcji wnętrza \footcite{idaPlatform2026}.

\section{Preferencje jawne (explicit)}
\label{ch03:sec:explicit}

Preferencje jawne obejmują wybory, które użytkownik potrafi zwerbalizować i świadomie wskazać: nazwy stylów, palety barwne, preferowane materiały oraz pożądany poziom jasności czy ciepła.
Charles Osgood, George Suci i Percy Tannenbaum zaproponowali technikę dyferencjału semantycznego do pomiaru konotacyjnych znaczeń pojęć na skalach rozpiętych między parami przymiotników przeciwstawnych \footcite{osgood1957measurement}.
W interfejsie IDA adaptacja ta przybiera postać intuicyjnych suwaków (m.in.\ ciepło, jasność, złożoność wizualna) oraz kafelków wyboru palet materiałowych.

Samoświadomość estetyczna bywa jednak zwodnicza. Deklaracja „preferuję chłodny minimalizm” nierzadko współwystępuje z entuzjastycznym wyborem bogatych fakturowo, ciepłych wnętrz w teście reakcji wizualnej. Dlatego obok deklaracji słownych niezbędny jest równoległy kanał preferencji ukrytych (rozdz.~\ref{ch03:sec:implicit}), a badanie rozbieżności między nimi stanowi jedno z głównych pytań badawczych pracy (\RQ{2}).

Do rejestru jawnego należy również przesłanie własnych obrazów inspiracji (rozdz.~\ref{ch03:sec:inspiracje}). Choć inspiracja operuje językiem formy, to użytkownik decyduje w pełni świadomie, \emph{jaki} obraz przedstawić systemowi jako wzorzec.

\section{Preferencje ukryte (implicit) i afektywna selekcja kart}
\label{ch03:sec:implicit}

Anthony Greenwald i współpracownicy wprowadzili test utajonych skojarzeń (Implicit Association Test, IAT), wykazując, że pomiar czasów reakcji na bodźce pozwala badać postawy i preferencje z pominięciem filtrów racjonalizacji i autoprezentacji \footcite{greenwald1998iat}.
W badaniach nad interakcją człowieka z komputerem (HCI) oraz w systemach rekomendacyjnych adaptacją tej idei stała się technika szybkiej, binarnej selekcji kart wizualnych (ang.\ \emph{card swiping}). Szybki gest odrzucenia lub polubienia obrazu pozwala na uchwycenie pierwotnej, afektywnej reakcji estetycznej.

W platformie IDA kanał ten przybrał formę modułu afektywnej selekcji kart: uczestnikowi prezentowana jest seria starannie dobranych wizualizacji wnętrz, które akceptuje lub odrzuca. System rejestruje kierunek wyboru oraz czas decyzji (\texttt{reaction\_time\_ms}), budując zagregowany profil preferencji intuicyjnych zasilający wariant ukryty (\texttt{implicit}) w macierzy.
Termin „ukryte” odnosi się w niniejszej pracy do pośredniej metody pomiaru behawioralnego i nie niesie psychoanalitycznych założeń o nieświadomości.

\begin{figure}[htbp]
  \centering
  \Todo{wstaw plik: figures/ch03/ch03-tinder-iat-flow.pdf --- schemat przepływu: bodziec wizualny $\rightarrow$ gest akceptacji/odrzucenia $\rightarrow$ czas reakcji $\rightarrow$ profil preferencji ukrytych}
  \caption[Schemat modułu afektywnej selekcji kart wizualnych w IDA. Opracowanie własne; rama teoretyczna IAT: Greenwald, McGhee i Schwartz; artefakt IDA.]{Schemat modułu afektywnej selekcji kart wizualnych w IDA.
    Opracowanie własne; rama teoretyczna IAT: \textcite{greenwald1998iat}; artefakt: \footcite{idaPlatform2026}.}
  \label{fig:ch03:tinder-iat}
\end{figure}

\section{Rozbieżność między deklaracją a zachowaniem}
\label{ch03:sec:rozbieznosci}

Daniel Kahneman opisuje architekturę poznawczą człowieka jako współdziałanie dwóch trybów: szybkiego, automatycznego i emocjonalnego Systemu~1 oraz wolnego, analitycznego i racjonalizującego Systemu~2 \footcite{kahneman2011thinking}.
W doświadczeniu estetycznym wnętrza oba te systemy często generują sprzeczne sygnały. \textcolor{orange}{Deklaracja słowna (System~2) bywa zdominowana przez gust aspiracyjny, mody kulturowe i potrzebę aprobaty społecznej, podczas gdy natychmiastowy gest afektywny (System~1) odsłania bezpośrednie zapotrzebowanie zmysłowe na poczucie przytulności, ciepła czy intymnego azylu.}

W platformie IDA stopień zgodności obu kanałów jest precyzyjnie kalkulowany (m.in.\ zgodność stylu, palety barw i nasycenia biofilią).
Pytanie~\RQ{2} weryfikuje, czy owa rozbieżność znajduje odzwierciedlenie w zachowaniu uczestnika przed macierzą wygenerowanych koncepcji: czy wydłuża czas namysłu nad wyborem oraz czy skłania do częstszego wybierania wariantów zrównoważonych (mieszanych).
Rozbieżność ta nie jest więc traktowana jako błąd pomiaru, lecz jako fundamentalna dramaturgia samopoznania estetycznego użytkownika.

\section{Wskaźniki behawioralne i czas wyboru}
\label{ch03:sec:behavioral}

Obok kierunku gestu system rejestruje czas reakcji uczestnika na poszczególne obrazy wnętrz. Czas ten przybliża stopień pewności decyzji: natychmiastowe odrzucenie sygnalizuje dysonans estetyczny, podczas gdy dłuższy namysł może świadczyć o ambiwalencji lub głębszej fascynacji kompozycją.

W ślepej macierzy kluczowym parametrem staje się łączny czas porównywania wariantów (\texttt{selection\_time\_ms}). W odróżnieniu od laboratoryjnego śledzenia mikroruchów oka, badanie terenowe IDA koncentruje się na naturalnym rytmie interakcji, chroniąc poczucie swobody uczestnika.

\section{Drabinowanie wartości (means-end laddering)}
\label{ch03:sec:laddering}

Thomas Reynolds i Jonathan Gutman sformułowali metodę drabinowania w teorii środków i celów (\emph{means-end chain}), ukazując, w jaki sposób konkretne atrybuty fizyczne obiektu łączą się z konsekwencjami funkcjonalnymi, by ostatecznie zakorzenić się w fundamentalnych wartościach egzystencjalnych człowieka \footcite{reynolds1988laddering}.
W module konfiguracji pokoju IDA dialog z awatarem prowadzi użytkownika po szczeblach tej drabiny: od pytania o konkretne wyposażenie i bariery użytkowe (punkty tarcia), przez ich wpływ na codzienne samopoczucie, aż po głęboką odpowiedź na pytanie, dlaczego dana zmiana jest dla niego ważna.

Dzięki temu generatywny brief przestaje być powierzchownym zbiorem przymiotników stylistycznych. Wiedza o tym, że poszukiwane „ciepłe drewno” ma sprzyjać wyciszeniu po stresującej pracy i budować poczucie bezpieczeństwa rodziny, pozwala na wygenerowanie wariantu funkcjonalnego (\texttt{mixed\_functional}) o autentycznym ładunku tożsamościowym.

\section{Inspiracje wizualne jako osobisty atlas wyobraźni}
\label{ch03:sec:inspiracje}

Przesłanie fotografii inspiracji wprowadza do dialogu preferencję przez konkretny okaz zmysłowy. Użytkownik wskazuje atmosferę, której złożoności nie potrafi wyrazić za pomocą suwaków czy słów.
W potoku IDA przesłane obrazy referencyjne są analizowane przez multimodalny model wizyjno-językowy, który wyodrębnia dominanty kompozycyjne, oświetleniowe i materiałowe. Wyniki tej analizy zasilają wariant referencji inspiracji (\texttt{inspiration\_reference}) w macierzy.

\section{Synteza wieloźródłowa i architektura wariantów}
\label{ch03:sec:synteza}

Wieloźródłowy model preferencji w IDA nie rości sobie pretensji do arbitralnego wskazywania „jedynego prawdziwego gustu”. Traktuje wrażliwość estetyczną jako zjawisko polifoniczne, przekładając poszczególne kanały na sześć zróżnicowanych wariantów koncepcji (tab.~\ref{tab:ch03:sources} i~\ref{tab:ch03:instruments}).

\begin{table}[htbp]
  \centering
  \caption{Sześć źródeł wytycznych generacji koncepcji wnętrza w eksperymencie IDA.}
  \label{tab:ch03:sources}
  \begin{tabular}{@{}llp{0.44\textwidth}@{}}
    \toprule
    Klucz źródła & Nazwa wariantu & Źródła danych wejściowych \\
    \midrule
    \texttt{implicit}
      & Wariant preferencji ukrytych
      & Gesty afektywnej selekcji kart, czasy reakcji \\
    \texttt{explicit}
      & Wariant preferencji jawnych
      & Dyferencjał semantyczny, deklarowana paleta i styl \\
    \texttt{personality}
      & Wariant osobowościowy
      & Aspekty cech inwentarza IPIP-NEO-120 (rozdz.~\ref{ch05:top}) \\
    \texttt{mixed}
      & Wariant zrównoważony (mieszany)
      & Fuzja danych ukrytych, jawnych i osobowościowych \\
    \texttt{mixed\_functional}
      & Wariant mieszany z funkcją pokoju
      & Fuzja preferencji + luka nastroju PRS, aktywności i wartości \\
    \texttt{inspiration\_reference}
      & Wariant referencji inspiracji
      & Analiza wizualna fotografii referencyjnej użytkownika \\
    \bottomrule
  \end{tabular}
\end{table}

\begin{table}[htbp]
  \centering
  \caption{Zestawienie instrumentów badania preferencji zintegrowanych w platformie IDA.}
  \label{tab:ch03:instruments}
  \begin{tabular}{@{}p{0.20\textwidth}p{0.22\textwidth}p{0.24\textwidth}p{0.24\textwidth}@{}}
    \toprule
    Instrument & Podstawa teoretyczna & Forma interakcji w IDA & Rola w procesie \\
    \midrule
    Mapa nastroju PRS
      & \textcite{hartig1997prs,kaplan1995restorative}
      & Siatka 2D nastroju
      & Pomiar luki regeneracyjnej \\
    Biofilia
      & \textcite{kellert2008biophilic}
      & 4 stopnie nasycenia
      & Warunkowanie obecności natury \\
    Dyferencjał semantyczny
      & \textcite{osgood1957measurement}
      & Suwaki atrybutów stylu
      & Pomiar preferencji jawnych \\
    Afektywna selekcja kart
      & \textcite{greenwald1998iat}
      & Gesty akceptacji/odrzucenia
      & Pomiar preferencji ukrytych \\
    Drabinowanie wartości
      & \textcite{reynolds1988laddering}
      & Dialog w Room Setup
      & Kontekst funkcjonalno-życiowy \\
    Inspiracje wizualne
      & Praktyka projektowa
      & Upload zdjęć referencyjnych
      & Bezpośredni transfer nastroju \\
    Big Five / IPIP-NEO-120
      & \textcite{john2008bigfive,gosling2002room}
      & Skale podwymiarów cech
      & Hipoteza osobowościowa \\
    \bottomrule
  \end{tabular}
\end{table}

\begin{figure}[htbp]
  \centering
  \Todo{wstaw plik: figures/ch03/ch03-multisource-model.pdf --- diagram: kanały preferencji $\rightarrow$ synteza briefu $\rightarrow$ macierz 6 wariantów $\rightarrow$ ślepy wybór}
  \caption[Wieloźródłowy model preferencji estetycznych w platformie IDA. Opracowanie własne na podstawie artefaktu.]{Wieloźródłowy model preferencji estetycznych w platformie IDA.
    Opracowanie własne na podstawie \footcite{idaPlatform2026}.}
  \label{fig:ch03:multisource}
\end{figure}

Sześć wygenerowanych koncepcji spotyka się w ślepym porównaniu. Użytkownik nie wie, które źródło wygenerowało dany obraz, oceniając wyłącznie siłę oddziaływania plastycznego i emocjonalnego przedstawionej przestrzeni. To właśnie ten moment stanowi empiryczny rdzeń pracy, weryfikowany w pytaniu~\RQ{3}.
